% =====================================================================
%  A1 -- Machine learning: linear regression by GD, feature scaling
% =====================================================================
\subsection{A1: Machine learning, regression and feature scaling}
\setprob{A1}

\begin{frame}{\probtitle{A1}{Linear regression by GD, and why feature scaling matters}}
  \begin{probbox}[Problem A1 \hfill Application: Machine Learning \quad $\star\star$]
    A model $\hat y = wx + b$ is trained on the data $(1,2),\,(2,3),\,(3,5)$ by minimising the
    mean-squared-error loss
    \[
      \mathcal L(w,b) = \frac{1}{2n}\sum_{i=1}^{n}\bigl(wx_i + b - y_i\bigr)^2, \qquad n = 3 .
    \]
    \vspace{-8pt}
    \begin{enumerate}[(a)]
      \item Derive $\nabla\mathcal L$ and take two GD steps from $(0,0)$ with $\alpha = 0.1$.
      \item Find the Hessian, its eigenvalues, $\kappa$, and the largest stable $\alpha$.
      \item Replace $x$ by the centred feature $\tilde x = x - \bar x$. Recompute the Hessian and $\kappa$, and
            compare the iterations needed for $10^{-6}$ accuracy at the best fixed step.
      \item Find the optimal $(w,b)$ and explain what you see.
    \end{enumerate}
  \end{probbox}
  \footnotesize\textcolor{mGrey}{The point: training \emph{is} gradient descent, and preprocessing is preconditioning.}
\end{frame}

\begin{frame}{\soltitle{A1}{1}{4}}
  \textbf{(a) Gradient.} With residuals $r_i = wx_i + b - y_i$,
  \[
    \frac{\partial\mathcal L}{\partial w} = \frac1n\sum_i r_i x_i, \qquad
    \frac{\partial\mathcal L}{\partial b} = \frac1n\sum_i r_i .
  \]
  \textbf{Step 1} from $(0,0)$: $\;r = (-2,-3,-5)$ and $\;\mathcal L = 6.3333$.
  \[
    \nabla\mathcal L = \Bigl(\tfrac{-2-6-15}{3},\ \tfrac{-10}{3}\Bigr) = (-7.6667,\,-3.3333)
    \;\Rightarrow\; (w_1,b_1) = \ans{(0.7667,\ 0.3333)}
  \]
  and the loss drops to $\mathcal L(w_1,b_1) = 1.2826$.
  \textbf{Step 2}: the predictions are $(1.1000,\,1.8667,\,2.6333)$, so $r = (-0.9000,\,-1.1333,\,-2.3667)$.
  \[
    \nabla\mathcal L = \Bigl(\tfrac{-0.9-2.2667-7.1}{3},\ \tfrac{-4.4}{3}\Bigr) = (-3.4222,\,-1.4667)
  \]
  \[
    (w_2,b_2) = \ans{(1.1089,\ 0.4800)}, \qquad \mathcal L = 0.2807 .
  \]
\end{frame}

\begin{frame}{\soltitle{A1}{2}{4}}
  \textbf{(b) Hessian.} $\mathcal L$ is quadratic, so the Hessian is constant:
  $\;\bH = \frac1n\sum_i\begin{pmatrix}x_i^2 & x_i\\ x_i & 1\end{pmatrix}
        = \begin{pmatrix}14/3 & 2\\ 2 & 1\end{pmatrix}.$
  \smallskip
  $\lambda^2 - \tfrac{17}{3}\lambda + \tfrac23 = 0 \;\Rightarrow\; \lambda = \dfrac{17\pm\sqrt{265}}{6}$, so
  \[
    \lmax = 5.5465,\quad \lmin = 0.1202,\quad
    \ans{\kappa \approx 46.1},\quad \ans{\alpha < \tfrac{2}{\lmax} \approx 0.3606}
  \]
  The best fixed step is $\alpha^\star = \tfrac{2}{\lmin+\lmax} = \tfrac{6}{17} \approx 0.353$. Its rate is
  $\tfrac{\kappa-1}{\kappa+1} = 0.9576$, which means \alert{319 iterations} to reach $10^{-6}$. For three data points!

  \begin{keybox}[Why is such a tiny problem so badly conditioned?]
    Every $x_i$ is positive ($\bar x = 2$), so the columns $[\,x,\ \mathbf 1\,]$ are strongly correlated.
    Nudging $w$ or nudging $b$ changes the fit in almost the same way, and that makes a long, narrow valley.
  \end{keybox}
\end{frame}

\begin{frame}{\soltitle{A1}{3}{4}}
  \textbf{(c) Centring.} $\tilde x = x - 2 = (-1,0,1)$, so $\overline{\tilde x} = 0$ and $\overline{\tilde x^2} = \tfrac23$:
  \[
    \tilde\bH = \begin{pmatrix}2/3 & 0\\ 0 & 1\end{pmatrix}
    \;\Rightarrow\; \ans{\kappa = 1.5},\quad
    \alpha^\star = \frac{2}{2/3+1} = 1.2,\quad \rho = \frac{0.5}{2.5} = 0.2 .
  \]
  Iterations for $10^{-6}$: $\;\lceil \ln10^6/\ln 5\rceil = \ans{9}$ instead of $319$. One subtraction made it \alert{35 times faster}.

  \centering
  \includegraphics[width=0.54\textwidth]{fig_scaling}
\end{frame}

\begin{frame}{\soltitle{A1}{4}{4}}
  \textbf{(d) The optimum.} Setting $\nabla\mathcal L = \mathbf 0$ gives the normal equations
  (using $\sum x_i = 6$, $\sum x_i^2 = 14$, $\sum y_i = 10$, $\sum x_iy_i = 23$):
  \[
    14w + 6b = 23,\qquad 6w + 3b = 10
    \;\Longrightarrow\; \ans{w = 1.5,\;\; b = \tfrac13}, \qquad \mathcal L^\star = \tfrac1{36} \approx 0.0278 .
  \]
  The centred model gives $\hat y = 1.5(x-2) + \tfrac{10}{3}$, which is \emph{the same line}. Centring changes
  the shape of the optimisation problem, not its answer.

  \vspace{4pt}
  \begin{keybox}[The bigger picture]
    \small
    \begin{itemize}
      \item Centring and standardising features is a change of variables that acts as a \emph{preconditioner}:
            GD sees nearly round contours instead of a narrow valley.
      \item Here a closed form exists. With millions of parameters there is none, GD or SGD is the only practical
            solver, and conditioning decides how long training takes.
    \end{itemize}
  \end{keybox}
\end{frame}
